% R152 source proposal; provenance and exact prior text remain in the proposal manifest.
\section{From a reach to a testable model}
Start with a movement we can measure well. For the proposed snap showcase,
the arm begins extended with the palm down. The elbow and forearm swing and
turn; the wrist and hand turn upward during the snap, then move down to rest.
The model must account for that whole movement, including the preparation,
contact and return. Finger pairing, organism-specific anatomy and timing still
need a real-reference series; a human-looking animation cannot establish that
a particular monkey can perform the same gesture.

The cup reach supplies a simpler precision comparison: relate recorded signals
to measured movement before attempting the coordinated snap. Each task still
needs its own anatomy, contact, timing and independent comparison.

That comparison connects three interacting groups: task-linked brain
observations, named-muscle observations and a musculoskeletal simulation.
A learner would relate the observations to the mechanics; independent movement
and contact measurements would test its predictions. EEG records scalp
electrical activity and surface EMG records muscle electrical activity.
A separately recorded nerve or spinal signal needs its own declared channel.
The proposed apparatus uses two arm sleeves, scalp EEG and synchronized,
calibrated left and right tripod views. Vendor, channel, bandwidth and clock
qualification remain open.

Intent to action is the target. By understanding and measuring the intermediary,
we aim to make that mapping as close to one-to-one as possible within a declared
task. Specify what counts as the intended action, the participant and session,
state and context, permitted history, tolerated error and independent reference.
The test then asks whether observing the middle reduces ambiguity and improves
prediction or qualified control. It does not assume a unique hidden intention
or claim that the best possible mapping has already been achieved.

The proposed learner must explain the recorded brain and muscle activity,
where the hand moved and when. These are proposed comparisons: no paired trial,
trained model, data acquisition, fine-tuning, reinforcement-learning rollout or
accepted checkpoint is established for this EEG/EMG/movement loop. Its sources
and permissions, physical and computational behaviour, clinical interfaces,
reference comparisons and held-out tests still need independent assessment.

The completed LINK comparison supplies a warning about what would count as
success. Its 18 runs and 720 saved epochs tested a different question; both
model-admission gates and all five benefit claims failed. Neither those runs
nor a future successful prediction would show that a trainee received sensation
or produced an action. Recording, prediction, simulated movement, sensory
delivery and physical action each need evidence of their own. EEG and EMG
recorders alone cannot write sensation, and no sensory delivery or physical
actuation is established for this proposal.

\subsection{Form, function and action need a measured middle}
A structural map constrains what a system could do. A NeuroConnectome maps
specified neural connections; its organism, coverage, resolution and source
would have to be declared. Following the reach through that map also requires
the signals and the processes that turn those constraints into an outcome.
Signals must remain identified as observed or inferred. This is why the proposed
programme follows form to function to action: understanding and measuring the
middle could improve the intent-to-action mapping for the declared task.

The molecular comparison asks the same question about a different mechanism.
DNA to RNA and RNA to protein put a process between a starting constraint and its
downstream product. Regulatory context, RNA processing and cell state change
what that process produces; protein structure alone does not determine the
behaviour of a whole cell. AlphaGenome, RNA or cell-state models and AlphaFold
illustrate these different modelling jobs. AlphaTranscriptome names a desired
predictive capability here. Connecting the predictors would still leave system
dynamics and the environment to account for, with experimental outcomes needed
to test the combined prediction.
% Source-role pointers: alphagenome2026, scgpt, alphafold3; bounded owner checks.
\par\noindent \href{https://www.nature.com/articles/s41586-025-10014-0}{AlphaGenome regulatory prediction}~\cite{alphagenomeRetained};
\href{https://www.nature.com/articles/s41592-024-02201-0}{scGPT cell-state modelling}~\cite{scgptRetained};
\href{https://www.nature.com/articles/s41586-024-07487-w}{AlphaFold 3 structure prediction}~\cite{abramson2024}.
These examples support the comparison of modelling roles only. The complete
AlphaGenome and scGPT methods have not been independently assessed here, and
the examples provide no evidence of an integrated cell programmer.

For the reach, the corresponding biological question is what happens between
brain activity and the body's response. Brain2Spine and Spine2Body name the
proposed links to test; their composition is the Brain2Body aim. The spinal cord
integrates descending and sensory signals and contributes local coordination,
reflex regulation and adaptation. RNA processing is a different mechanism;
the shared question is how a context-dependent process changes an initial
constraint into an outcome. Supported spinal observations would test that
account. Where they are absent, internal states remain inferred or simulated.
Physical spinal work remains read-only.
\par\noindent \href{https://pubmed.ncbi.nlm.nih.gov/38603479/}{Spinal sensorimotor adaptation}~\cite{lavaud2024}
offers a constrained example of that processing. It validates neither our
proposed pathway nor unrestricted skill transfer. We draw on the abstract and
relevant author-manuscript passages; equivalence to the final published edition
and the complete methods and statistics remain unverified here.

Follow the action one step further and the question changes: what did it do
to the task environment? Thought2Action relates a declared intention to an
independently measured action. Action2World follows the action to its consequence
in that environment; Thought2World names the proposed composition. The earlier
Body2World label drew attention to the consequence outside the body. Body2Action
focuses on body mechanics and a specified task outcome, while Brain2Motion
focuses on the movement pathway. Each working label identifies a related scope
for a test; the labels establish neither interchangeable nor working controllers.

The engineering route would bring these observations and comparisons into a
learner: neural activity to musculoskeletal simulation to a programmable model,
toward NeuroTranslation. The proposed movement-focused learner is a large
movement model. It would perform the declared prediction or control. System 3
thinking names the proposed human--AI supervisory partnership around it: a
supervisory model would compare task strategies, diagnose where neural or
mechanical predictions miss independent measurements, and propose training or
calibration comparisons for the signal mapping and virtual body. The learner
would answer the task; the supervisor would help choose what to test next.
Neither replaces the biological spinal intermediary or a physical interface.
Final held-out outcomes must remain outside rewards and model selection.
\par\noindent Shaw and Nave's SSRN preprint,
\emph{Thinking---Fast, Slow, and Artificial}, and its
\href{https://executiveeducation.wharton.upenn.edu/thought-leadership/wharton-at-work/2026/05/thinking-fast-slow-and-artificially/}{institutional discussion}
provide conceptual context for this proposed human--AI role.~\cite{shawNaveSystem3Retained}
We use that framing only; the complete preprint and institutional article have
not been reviewed here. Neither establishes a third anatomical brain system or
validates the proposed learner or supervisor. Expert footage may supply movement
and tactical examples, with rights and camera uncertainty accounted for; it
supplies no measured expert EEG, EMG, muscle force or stimulation plan.

The outcome then supplies the next question for the learner and its supervisor.
Action changes what happens; its sensory consequences inform perception and
the next intention. Independently measured movement would anchor the action
part of that loop, while known stimuli, measured feedback and appropriate
recordings or reports would have to test the further sensory mappings. Movement
alone would not decode a unique sensation or perception. The return path must
be measured rather than drawn into existence. The full Interoceptomics aim
includes internal-body sensing and connectivity; external touch and
proprioceptive tasks contribute evidence but do not by themselves solve that
wider problem. A prediction, assisted execution and retained skill without
assistance answer different questions. Each wider claim must earn its own
independent comparison.

\subsection{Context connects. Grounding checks}
Within a reach, context has a concrete job: relate brain activity to muscle
activity in this task and at this time. Attention is one proposed way to
represent those relationships. The model compares a representation called a query with representations called keys, then uses the resulting weights to combine their values. Masks keep unavailable observations and future
samples out. The calculation weights information; it does not establish that
attention understands all context.
% Source registry IDs: context-native-01a12f62, context-attention2017-v7, context-mult2019, context-policy4
\par\noindent \href{https://arxiv.org/html/1706.03762v7}{Vaswani et al. (2017)}~\cite{vaswani2017}; \href{https://aclanthology.org/P19-1656/}{Tsai et al. (2019)}~\cite{tsai2019}. These methods motivate the computation; validation of the proposed paired system still requires the independent observations below.

A pattern can fit inside the model and still describe the wrong movement.
Grounding therefore means testing the declared prediction against qualified
observations under the same conditions. Independently measured movement would
check the outcome, while muscle activity would add an observation that
constrains the brain model. Even a matching reach cannot identify a unique
private experience. The proposed paired trial would have to show which of
these observations adds useful information.

\subsection{Muscle activity tests the brain model}
Record brain and named-muscle activity during the same reach. In the proposed
paired trial, EEG would sample scalp electrical activity and surface EMG would
sample activity associated with the named muscles. Paired training could use
the muscle observation to constrain their relationship. Then withhold the
muscle target and movement reference from the EEG-only predictor: can the
permitted brain observations predict a specified muscle-activity observable
or movement?

Pairing and calibration remain proposed. Shared cues and movement artefacts
could give that predictor a shortcut, so the comparison needs controls for
both. The muscle observations would cover only the named muscles, and EMG
supplies neither afferent nerve activity nor direct force. Sensory claims
still need known stimuli and measured feedback. Several neural states could
remain compatible with even a well-matched pair; the pair could narrow a
task-specific explanation without revealing the complete hidden brain.
Checking the movement itself requires the next independent observation.
% Source registry IDs: context-native-01a12f62, context-policy4, overhaul-brantley2018
\par\noindent This proposed pairing still needs supporting evidence. The candidate reference at \url{https://www.nature.com/articles/sdata2018133} has not yet been checked for exact support or reuse rights.

\subsection{Measured motion tests the virtual body}
Measure the same hand and arm from the aligned left and right tripod views.
The proposed motion-capture relay-composite would align the two point maps
in a common clock and coordinate frame, then join them into one composite
while retaining each point's origin. Before that joined map can check a
prediction, its calibration, alignment and capture precision would have to
be independently verified.

The movement map would constrain the external movement model: surface motion
and pose in the task's physical frame. To explain that motion mechanically,
the internal musculoskeletal model would describe specified bones, joints
and soft tissues. Organism, developmental stage and task would determine the
appropriate anatomical sources. X-ray projections could inform bone and
joint relationships; appropriate MRI references could inform specified soft
tissues. Registration, rights and source-specific limits would need independent
checks. No scan set is selected for this proposed comparison. Anatomy still
leaves activation, forces and neural function to test.

With a qualified composite, the musculoskeletal-simulation physics-engine's
joints, muscles, contacts and physical limits would face a measured comparison.
Fit the model on training data, then compare its predicted reach with
held-out measured movement. A force claim would require an additional force
measurement. The body's form limits possible motion; measured function would
test the mechanics and show where they need refinement.
% Source registry IDs: context-native-01a12f62, context-component-names-01a12f81, context-policy4, context-opensim2018, context-myosuite-docs, context-precision-01a12f77
\par\noindent \href{https://journals.plos.org/ploscompbiol/article?id=10.1371/journal.pcbi.1006223}{Seth et al. (2018)}~\cite{opensim2018}; \href{https://myosuite.readthedocs.io/en/latest/}{MyoSuite documentation}~\cite{myosuiteDocsUndated}. The simulation tools supply methods for that comparison; the proposed capture and body models still require independent validation.

A lifelike hand could expose a missed contact or a late reach. The comparison
still needs independent capture alignment, precision, mechanics and biological
fidelity; none has been demonstrated for this proposal. Captured geometry
cannot establish hidden neural state, internal forces or unique subjective
experience. It would check the declared movement prediction within those limits.

\subsection{Two pairs. One conditional learner}
The learner would have two connected jobs: relate EEG observations to muscle
activity, and relate measured movement to predicted mechanics. Large movement
model names the proposed family; deep neural network describes a possible
implementation. Large refers to intended scope, not a demonstrated parameter
count. Encoders would accept only observations permitted for the named test.
Two branches would summarize those relationships in a shared task
representation, with separate output parts for each direction's targets.
% Source registry IDs: context-native-01a12f62, context-policy4, context-attention2017-v7, context-mult2019
\par\noindent \href{https://arxiv.org/html/1706.03762v7}{Vaswani et al. (2017)}~\cite{vaswani2017}; \href{https://aclanthology.org/P19-1656/}{Tsai et al. (2019)}~\cite{tsai2019}. These computational methods offer precedents for a design, leaving the proposed paired learner and its performance to be tested.

Signal features belong to branches, and branches belong to the learner.
That hierarchy describes the computation, not two anatomical thinking systems.
A common clock, task and physical frame would connect the branches; independent
measurements would check their errors. Form constrains the prediction, and
checked function would guide refinement. Several hidden states could still
fit the same reach, so neither a shared representation nor conditional
predictions would establish an exact inverse nervous system. The architecture
has not been selected, implemented, trained or independently accepted.
To choose it, we would have to compare it with simpler baselines on
appropriately reserved observations.

\subsection{Learn the movement, then test what transfers}

A recording of a reach already contains a possible lesson for a computer: use the earlier movement to predict where the hand goes next. It need not contain EEG to pose that question. A separate proposed comparison would ask whether learning from one person's grasp-and-lift movement helps predict another person's movement. The WAY-EEG-GAL description identifies continuous recordings of lifting series as well as separate windowed trials and lift-summary tables.\citep{luciw2014way} The continuous movement could supply both the history and the later checking observation. It would not require matching AllLifts Run or trial rows, and it would create no measured neural pair between the two people. The original study of five complete weight series remains a different design.

To transfer anything useful, we must say what corresponds. Homologous structures share an evolutionary anatomical origin. Mechanically analogous structures can perform a comparable job without that shared origin: a robot gripper's opening and a human thumb--finger opening might both describe how an object is approached. That example would require a separate robotic study; the proposed two-person study would not test it. Each body keeps its own joints and possible directions of movement, positive-axis conventions, proportions, muscle--tendon roles, loads, inertia, contact and timing. Two similar-looking limbs need not generate the same forces or move at the same rate. Calculating a corresponding movement for another body is retargeting; the calculation remains derived, alongside separately identified tracking observations, estimates, predictions and simulations.

The proposed description of movement contains thirteen quantities. Three give the index sensor's position relative to the wrist, three do the same for the thumb and three for the object. One describes the distance between the two digit sensors; the final three describe wrist velocity using only information already available. Thirteen extra mask values say which features are actually available. Each person needs checked position units, sensor placement and a wrist orientation that turns positions into the intended frame without reflecting the body. After the checked conversion to metres, the middle, or median, wrist-to-index-sensor distance from that person's assigned training support supplies a positive length scale for position and velocity. This is a sensor-distance scale, not a measured bone length; development observations do not set it. Velocity is also multiplied by the same 0.1-second interval in both people. The thirteen prediction features then have no physical units. To recover velocity errors in metres per second, multiply by that person's length scale and divide by 0.1 seconds. Missing features contribute no target error; writing a masked zero does not make a missing observation measured. These features do not supply all finger joints, individual muscle forces or the person's inertia.

The model would use fifty movement frames at 100 observations per second to predict 100 milliseconds beyond the latest filtered input actually available. The source rate must first be checked as 500 observations per second. A proposed filter combining 51 present and earlier samples would precede keeping every fifth sample, producing 100 observations per second. Its linear-phase delay is 50 milliseconds. That delay stays in the account of when the information became available; moving a timestamp backward cannot make the predictor respond earlier. Filter coefficients, the frequencies it preserves, timestamp conventions, gaps and handling at recording edges still need checking. Only the named movement columns may enter. The source's kinematics matrix also contains force-plate signals, which cannot quietly become movement features.

Knowing which input was read is different from knowing whether two inputs are independent. Trustworthy published participant attribution, an unchanged exact archive entry and a checked capture can identify the participant record and recorded input without requiring one particular spelling of an ID inside the file. They do not authenticate a biological person or prove that two files came from different original acquisitions. A longer recording and its trimmed copy could have different names and file fingerprints while sharing the same movement. Shared samples, trims, derived versions and overlaps must therefore remain together as one acquisition group before filtering, cutting windows, scaling or assigning roles. The complete public non-ST roster and these relationships still need establishing. All 294 rows of the earlier AllLifts table remain unavailable for training, evaluation and assigning trials or acquisition groups. ST, competition storage and reserved LINK recordings remain closed.

Qualified P1 groups would provide the earlier movement learning. P2 groups would be divided between target training and development checking. At least two independently established P2 acquisition groups are needed. If there are $G\geq2$, a fixed hash calculation ranks their established original-parent identifiers; the first $\lfloor3G/5\rfloor$ become training groups and the rest become development groups. The floor brackets mean rounding down, so the nominal 60/40 division is explicit even when the number of groups is small. The exact identifiers, how they are written as bytes and the complete roster must be settled before ranking or seeing outcomes. The order is not chronology and need not balance lifting conditions. If uncertain ancestry forces all P2 recordings into one group, the comparison stops. Cutting that group into windows cannot create independent training and checking acquisitions. No actual group or movement segment has been assigned by this proposed identity method.

The three training arms keep the target task and model the same. B0 first learns from correctly mapped P2 training movement. B1 and B2 first learn from exactly the same P1 windows in the same order. B2 uses the proposed correct correspondence. B1 deliberately uses the thumb sensor for the index role, the object sensor for the thumb role and the index sensor for the object role, leaving the wrist anchor fixed. Both use the same correct source calibration and scale before this change, then the same correct target adapter and ordered target-training batches. Sampling gives each original acquisition group balanced treatment and records repeated as well as unique windows. The learner might undo the false correspondence; changing roles can also change which features are available. Equal numbers of updates do not give the arms equal diversity of source movements. Those limits must stay in the explanation of any difference, rather than prompting a more convenient control after results are known.

The proposed network is a single gated recurrent unit, or GRU, layer: a calculation that updates a numerical memory as it reads the sequence. It has a memory width of 32, accepts the thirteen movement features and thirteen masks as 26 inputs, and predicts thirteen outputs. That software memory is not measured biological memory. Each arm uses the same initialization for a given fixed random seed, with seeds 17, 31 and 73. Every arm/seed receives 200 earlier-learning updates and 100 target-training updates, using batches of 16 examples. Adam, the proposed weight-update rule, has learning rate 0.001, beta settings 0.9 and 0.999, epsilon $10^{-8}$ and zero weight decay. There is no clipping, changing learning-rate schedule or early stopping; the final update supplies the checkpoint. During training, squared prediction errors are averaged over observed target features; a batch with no observed targets must stop. Across three arms and three seeds this would be 2,700 updates. None has been performed for this comparison. No external checkpoint is selected, and the numerical updates, filter, precision and actual computing limits still need independent checks.

For each development acquisition group and seed, the proposed score adds the squared errors for the observed target features, divides by how many such targets were observed, then takes the square root. It then averages those group scores equally, so a long recording does not automatically outweigh a short one. A comparison with no observed targets has no valid denominator and must stop. The report must give feature counts, missing observations, errors converted back to physical units, uncertainty and a persistence forecast---continuing the latest available movement value---using the same information endpoint. Every worse group or feature remains visible. B2 would need lower normalized error than both B0 and B1 in each fixed seed, with no applicable original requirement failing, to justify another investigation. Two people and development checking do not establish performance across a population, untouched final confirmation or biological teaching.

Larger movement collections could supply more examples after their exact files, permissions, relevant movements and original-recording overlap are checked. \href{https://www.addbiomechanics.org/download_data.html}{AddBiomechanics} distinguishes recorded markers and ground reactions from fitted joint movements and torques estimated by inverse dynamics---calculating forces or torques consistent with motion and a body model. Those estimates are not new direct force measurements. The selected \href{https://amass.is.tue.mpg.de/license.html}{AMASS} asset would need its research-use, non-commercial and redistribution terms checked. The same original recording in two collections stays one source, not independent training and checking data. LaBraM's earlier training inventory flags the WAY source family, but exact checkpoint, participant and trial exposure remains unknown. A later decision to reserve a recording cannot remove its possible influence from already trained weights. No such checkpoint is selected for this comparison.

Scans have a more specific job than making the body look complete. Before adding X-ray/CT or MRI, check which anatomical sources the chosen musculoskeletal model already uses. A new scan is useful when it fills a stated gap or provides a genuinely independent check. Its person or species, scale, registration, permissions and uncertainty still need checking. Registration establishes the correspondence between that image and the body being studied. Another person's scan remains an external reference; a generic body model does not secretly measure this participant's geometry, muscle activation or forces. No scan set is selected here.

The actual eligible recordings, acquisition groups, units, frames, clocks and calibration are still unresolved for this continuous comparison. It reports no new fit or successful transfer. A movement description that proved useful could later be tested alongside genuinely paired neural and body observations, using only the inputs allowed for that prediction. The transferred or calculated target would not become a simultaneous measurement. The useful next answer is where the movement forecast improves, where it worsens and which measurement or assumption explains the difference. Neural equivalence, causal physiology, delivered sensation and a skill that remains after guidance ends would each need a further experiment.

\subsection{Pretrain. Pair. Verify. Refine}
Those comparisons start with the observations we may use. Data rights,
calibration and held-out subjects, sessions and tasks determine which can
enter the study; resolve them before processing data or choosing a model.
Open-source code does not grant blanket permission to use datasets,
checkpoints or clinical interfaces. Qualified data could first train
representations for each recording type. LaBraM is one EEG encoder candidate,
not a complete brain-body writer. Synchronized paired observations could
then train the declared mapping from permitted inputs.

Paired prediction would give later refinement a baseline to beat. Before
optional reinforcement learning, post-training or task-specific fine-tuning,
the physical predictions and reward need checking against independent
outcomes; a reward the computer can calculate supplies no biological proof
by itself. Offline reinforcement learning needs qualified coverage of actions,
transitions and rewards: labels alone cannot tell us what an untried action
would do. Backpropagation can update differentiable neural components, but
it cannot automatically pass through every physics engine; universal
differentiability is not established. Simpler baselines and frozen final
tests remain through every phase.
% Source registry IDs: context-native-01a12f62, context-policy4, context-labram2405, context-cql2020-abstract, context-opensim2018, context-myosuite-docs
\par\noindent \href{https://arxiv.org/html/2405.18765v1}{Jiang, Zhao and Lu (2024)}~\cite{labram2024}; \href{https://arxiv.org/abs/2006.04779v3}{Kumar et al. (2020)}~\cite{cql2020}; \href{https://journals.plos.org/ploscompbiol/article?id=10.1371/journal.pcbi.1006223}{Seth et al. (2018)}~\cite{opensim2018}; \href{https://myosuite.readthedocs.io/en/latest/}{MyoSuite documentation}~\cite{myosuiteDocsUndated}. These sources offer candidate methods for the sequence; validation of the proposed paired system still requires the experiment below.

This proposed paired training sequence has not been carried out.

\subsection{Respond quickly. Improve carefully}
A frozen accepted model would provide the proposed response: predict from
permitted past and present observations, then prepare encoded outputs for a
qualified interface. Inspecting the errors would be a separate job. That path
would calibrate and propose a candidate update, to be independently evaluated
with a restorable checkpoint before it could replace the running model.
Updating an estimate of the current state therefore differs from changing
learned weights; examining an error would not silently change the predictor.

The candidate update would need more than a good-looking new score. Its
independent evaluation, a saved version that can be restored, evidence that
restoration works and the applicable authority would all precede replacement
of the running model. Until then the accepted weights stay fixed. This keeps
an attempt to improve tomorrow's model from quietly changing the one being
tested today.

The fast-and-slow comparison is a Kahneman-like analogy for this division
of work. It identifies no pair of biological systems and assigns no inherent
thinking speed to EEG or spatial data. Whether the frozen response is quick
enough would depend on the complete path from acquisition through filtering,
alignment, inference, output, feedback and rendering, including delay tails
and jitter. That end-to-end latency has not been measured. No accepted model,
real-time demonstration or running system that adapts under these rules is
established here, so the proposed separation supplies no latency or accuracy
guarantee.
% Source registry IDs: context-native-01a12f62, context-policy4

The complete timing test would include preparing the output and any later
sensory conversion or actuation, rather than stopping the clock when the
prediction leaves the model. Delay tails describe the unusually slow cases;
jitter describes variation in delay. Missed deadlines, uncertainty and the
effect of other healthy workloads would also need reporting. An average
inference time could hide the late response that matters in a catch.

\subsection{Write sensation. Enact action}
What reaches the trainee when the computer has made its prediction? A sensory
transducer would have to deliver a sensory output; a controller and actuator
would have to produce action in a virtual, robotic or later qualified biological
body. Input and output describe the computer boundary. Afferent and efferent
describe physiological direction relative to the central nervous system.

A display or haptic interface supplies sensory feedback, which differs from direct neural writing. Direct cortical writing could have a sensory role without being a peripheral afferent signal, just as a robot command is not an anatomical efferent nerve. EEG and EMG recorders do not write, and predicting an EEG-like pattern does not deliver sensation. Each proposed output needs an independently qualified interface, permission for its use and a measurement of what it delivers. No such physical interface, stimulation, acquisition or actuation is established by this proposal.

No methods evidence from the inaccessible clinical source supports this
proposed interface. Implanted precedents cannot establish that EEG or EMG recorders are equivalent to devices that deliver sensation or produce an action. Those physical additions remain unpriced and outside the acquisition and virtual-modelling scope.
% Source registry IDs: context-write-native-01a12f66, context-native-01a12f62, context-policy4

\subsection{Close the loop with evidence}
Return to the cup with two independent checks: withheld body observations
for the brain-to-body prediction, and the recorded neural target for the
conditional estimate from the body. Agreement between the models or a high
reward could still accompany a wrong answer. Neither establishes a unique
inverse brain or someone's private experience. A later sensory or action
interface would need a further measurement of what it delivered, under which
conditions, with what error and delay.

Approximation gets us towards plausibility and virtualization gets us to
precision. Independent measurements, repeatability, causal checks and stated
error bounds must earn that precision; the proposal supplies neither perfect
biology nor an automatic result. The original held-out and real-reference
controls would test where each prediction works and where it fails. The
retained-learning controls ask a further question: what remains when useful
guidance stops, with any necessary movement bridge kept constant?

Physical sensory writers, transducers and actuators remain unpriced additions
outside the acquisition and virtual-modelling costs. Those costs cannot tell us
the price of delivering sensation or action. No completed loop, training or
physical-writing result is claimed here. \emph{Catch it again when the help is
gone} develops that task and retention comparison. \emph{Build a brain--body
system we can test} asks what those results would have to establish before the
prototype could grow. The mechanism, task and prototype are stages of the same
proposed experiment.

\subsection{What the attention calculation actually does}

The proposed learner needs a way to relate permitted observations. Consider
one moment in a reach. A query represents what the model is looking for;
keys represent the observations it could compare with that query; values
carry information from those observations. The model learns these
representations. Their names describe a calculation, not a person inside the
machine who understands the task.

The technical edition writes that calculation as
\[\operatorname{Attention}(Q,K,V)=\operatorname{softmax}\!\left(QK^{\mathsf{T}}/\sqrt{d_k}+M\right)V.\]
Here Q holds the queries, K the keys and V the values. The superscript on K
means transpose: rearrange its rows and columns so the comparisons can be
calculated. Multiplying Q by that rearranged K gives scores for the
query--key matches. The common width of each query and key is $d_k$; dividing
by its square root scales those scores. The softmax operation turns the
allowed scores into weights that sum to one for a query with usable admitted
keys. Multiplying by V then combines the values using those weights.

M is the admitted mask. It controls which observations may contribute, so
an unavailable muscle target or a later sample cannot quietly provide the
answer. A mask is part of the declared test, not a decorative option after
the model has seen the data. Channels, recording type, task and time must all
retain their meanings through the calculation. Even a correctly calculated
weighted combination could predict the wrong reach: the independent
measurement still has to check it. This is a proposed method; no weights,
EEG encoder or gradient route have been selected or executed here.

\subsection{What each direction is allowed to know}

Use the same short labels as the technical method: B means EEG, E means
EMG, Y means independently measured movement, C means declared task,
stimulus and context, and P means the physics model with its identified
source. The letters keep track of information. They do not make the
quantities interchangeable.

An EEG-only test would use permitted past and present B, together with
declared context, to predict a specified E or Y target. Training may have
used qualified paired examples. At prediction time, however, the withheld
muscle and movement answers and future samples must stay outside the
predictor. Alignment, filtering and preparation of the data must respect
that separation too. Supplying a target through a less obvious route would
still supply the answer.

A body-only test would use permitted E, Y or both to estimate specified
recorded neural features. Its answer is conditional: given these observations
in this task, what recorded features are plausible, and with what uncertainty?
Several hidden neural states could fit the same movement. The reverse
estimate therefore needs its own recorded target; it cannot recover a unique
inverse brain or deliver the sensation it predicts.

A hybrid test would use additional permitted recording types and need its
own named route and evaluation. Its result could not be presented as EEG-only.
Inputs calculated by the simulator would remain simulated. They would not
replace independent observations or inherit biological validity from a
lifelike hand.

Two models could also learn to agree with one another. That agreement is
called cycle consistency when following one mapping and its reverse brings
the result back toward its starting point. Both mappings could still be
wrong about the measured body. Shared cues, movement artifacts, subject
identity, overlapping sessions and information leaked during preprocessing
need the original controls. Final held-out answers must remain outside
selection of the model, physics and reward.

\subsection{How the proposed training would be checked}

First decide which observations may enter the study. Establish the rights,
people and tasks, physical units, calibration, uncertainty and shared clocks.
Freeze the relevant participant, session and task splits, then check whether
earlier pretraining already included material that would compromise those
splits. Those decisions precede preprocessing and model selection.

Qualified recordings could then train a representation for each recording
type. In the candidate LaBraM route, a neural tokenizer turns signal segments
into codes; masked-code prediction trains a model to estimate hidden codes;
later fine-tuning adapts it to a declared task. These are computer-learning
steps. LaBraM is a candidate EEG encoder, not a complete brain--body writer.
Permission to use code, checkpoints, datasets and individual recordings
requires separate qualification.

Next, paired supervised learning would use named measured targets and
explicit masks. Supervised means the training examples provide the answers
against which predictions are compared. Simpler predictors, cue-only
predictors and artifact controls would help test whether the proposed neural
route contributes useful information. Raw recordings, processed features and
model-generated values would remain distinguishable.

Combining errors requires another decision. An illustrative training
objective could include errors in EMG, movement and recorded neural targets,
applicable physical constraints and optional consistency terms. Before
choosing their weights, each term needs its target definition, units and
measurement uncertainty. Otherwise changing a unit or weight could change
which error the optimizer mainly tries to reduce. Consistency asks whether
the outputs fit together; an independent measurement asks whether they fit
what happened.

Optional reinforcement learning and post-training would come only after
checking the physical predictions and reward against independent outcomes.
Success on the task, smaller calibrated movement error and valid physical
states are possible reward ingredients, not results already achieved.
Loss and reward weights would be declared from training and development
evidence. Offline reinforcement learning needs observations that support
the relevant actions, resulting changes and rewards. A labelled trajectory
alone cannot tell us what an unsupported, untried action would do.
Conservative Q-Learning is one candidate method for this stage; its name
provides no safety or performance guarantee here.

Task-specific fine-tuning would still face frozen final tests in both
directions, with uncertainty, baselines and every original threshold.
Failure would remain a result. An accepted version could then run with its
weights fixed; a later update would need separate independent acceptance,
restoration proof and applicable authority before replacing it.

Backpropagation calculates how to change differentiable parts of the
computer model. It is different from biological feedback. A physics engine
does not automatically supply the required derivatives: the actual
OpenSim/MyoSuite simulation and optimization route, or a separately validated
substitute, would need checking. No such training, reinforcement-learning
rollout or task checkpoint is created or certified here. The procedure makes
the proposed comparison inspectable while leaving its outcome open.
% Source registry IDs: context-native-01a12f62, context-write-native-01a12f66, context-precision-01a12f77, context-policy4, context-spec15, context-glossary13, b2b-circular-checklist-original
